\chapter{عضویت، تجزیه و الگوریتم \lr{CYK}}

وقتی یک گرامر داریم، طبیعی‌ترین سؤال این است: «آیا رشته $w$ در زبان این گرامر هست؟» و اگر هست، «چطور تولید می‌شود؟» کامپایلر دقیقاً همین کار را با کد برنامه انجام می‌دهد. در این فصل چند الگوریتم برای این دو سؤال می‌بینیم.

\section{صورت جمله‌ای، عضویت و تجزیه}

\begin{definition}[صورت جمله‌ای]
اگر $w\in L(G)$ و $S\Rightarrow w_1\Rightarrow\cdots\Rightarrow w_n\Rightarrow w$ یک اشتقاق برای $w$ باشد، به $w$ یک \textbf{جمله} و به رشته‌های میانی $S,w_1,\dots,w_n$ که هم متغیر و هم پایانه دارند \term{صورت جمله‌ای}{Sentential Form} می‌گوییم.
\end{definition}

\begin{example}
در گرامر $S\to aSb\mid\eps$ داریم $S\Rightarrow aSb\Rightarrow aaSbb\Rightarrow aabb$، پس $S\der aabb$. رشته $aabb$ یک جمله از زبان است و $aaSbb$ یک صورت جمله‌ای.
\end{example}

برای گرامر $G$ و رشته $w$ دو مسئله داریم:
\begin{itemize}
\item \textbf{عضویت:} آیا $w\in L(G)$؟ به الگوریتمی که این را مشخص کند \textbf{الگوریتم عضویت} می‌گوییم.
\item \textbf{تجزیه} (\lr{parsing}): اگر $w\in L(G)$، دنباله قاعده‌هایی را پیدا کن که $w$ را تولید می‌کنند.
\end{itemize}

\section{تجزیه با جست‌وجوی کامل}

ساده‌ترین روش، امتحان کردن همه اشتقاق‌هاست (\textbf{تجزیه با جست‌وجوی کامل} یا \textbf{brute force}): در دور اول همه قاعده‌های $S\to x$ را به کار می‌بریم. در هر دور بعدی، چپ‌ترین متغیر هر صورت جمله‌ای را با همه قاعده‌های ممکن جایگزین می‌کنیم. اگر $w\in L(G)$، اشتقاق چپ‌ترینی با طول متناهی دارد، پس این روش بالاخره آن را پیدا می‌کند.

\begin{example}
$w=aabb$ و گرامر $S\to SS\mid aSb\mid bSa\mid\eps$.
\textbf{دور اول:} $S\Rightarrow SS$، $S\Rightarrow aSb$، $S\Rightarrow bSa$ و $S\Rightarrow\eps$. دو صورت آخر به درد نمی‌خورند ($w$ با $b$ شروع نمی‌شود و $\eps\neq w$)، پس کنار می‌روند.
\textbf{دور دوم:}
\[\begin{array}{llll}
S\Rightarrow SS\Rightarrow SSS, & S\Rightarrow SS\Rightarrow aSbS, & S\Rightarrow SS\Rightarrow bSaS, & S\Rightarrow SS\Rightarrow S,\\
S\Rightarrow aSb\Rightarrow aSSb, & S\Rightarrow aSb\Rightarrow aaSbb, & S\Rightarrow aSb\Rightarrow abSab, & S\Rightarrow aSb\Rightarrow ab.
\end{array}\]
در دور بعد از $aaSbb$ با $S\to\eps$ به $w$ می‌رسیم.
\end{example}

این روش دو مشکل جدی دارد: \textbf{ناکارآمد} است و ممکن است برای رشته‌ای که در زبان نیست (مثلاً $w=abb$) \textbf{هرگز متوقف نشود}. مشکل توقف از قاعده‌های $A\to\eps$ و $A\to B$ می‌آید، چون با آن‌ها طول صورت جمله‌ای می‌تواند کم شود یا ثابت بماند.

\begin{example}
همان زبان (بدون رشته تهی) را گرامر زیر بدون قاعده $\eps$ و قاعده واحد تولید می‌کند:
\[S\to SS\mid aSb\mid bSa\mid ab\mid ba.\]
برای هر $w\in\{a,b\}^+$ جست‌وجوی کامل حداکثر بعد از $\abs{w}$ دور تمام می‌شود، چون در هر دور طول صورت جمله‌ای دست‌کم یکی زیاد می‌شود. بعد از $\abs{w}$ دور یا اشتقاقی پیدا شده یا می‌دانیم $w\notin L(G)$.
\end{example}

\begin{theorem}
فرض کنید $G=(V,T,S,P)$ گرامر مستقل از متنی بدون قاعده‌های $A\to\eps$ و $A\to B$ باشد. آنگاه جست‌وجوی کامل را می‌توان به الگوریتمی تبدیل کرد که برای هر $w\in T^*$ یا اشتقاقی برای $w$ پیدا می‌کند یا اعلام می‌کند که اشتقاقی وجود ندارد.
\end{theorem}
\begin{proof}
در هر گام اشتقاق دست‌کم یکی از این دو اتفاق می‌افتد: طول صورت جمله‌ای زیاد می‌شود (قاعده با سمت راست بلندتر از ۱) یا تعداد پایانه‌ها زیاد می‌شود (قاعده $A\to a$). هیچ‌کدام از این دو مقدار هرگز کم نمی‌شود، چون قاعده $\eps$ نداریم. همچنین در اشتقاق $w$، نه طول صورت جمله‌ای و نه تعداد پایانه‌ها نمی‌تواند از $\abs{w}$ بیشتر شود. پس اشتقاق $w$ حداکثر $2\abs{w}$ گام دارد. الگوریتم همه اشتقاق‌های چپ‌ترین با حداکثر $2\abs{w}$ گام را امتحان می‌کند و اگر هیچ‌کدام $w$ را نساخت، جواب «نه» می‌دهد.
\end{proof}

\textbf{هزینه:} اگر فقط اشتقاق‌های چپ‌ترین را بشماریم، بعد از دور اول حداکثر $\abs{P}$ صورت جمله‌ای داریم، بعد از دور دوم حداکثر $\abs{P}^2$ و همین‌طور تا آخر. پس تعداد کل صورت‌های جمله‌ای حداکثر
\[M=\abs{P}+\abs{P}^2+\cdots+\abs{P}^{2\abs{w}}=O\big(\abs{P}^{2\abs{w}+1}\big)\]
است که با طول رشته \textbf{به‌طور نمایی} رشد می‌کند. در عمل هم این روش بسیار کند است.

\begin{theorem}
برای هر گرامر مستقل از متن، الگوریتمی وجود دارد که هر $w\in L(G)$ را در تعداد گامی متناسب با $\abs{w}^3$ تجزیه می‌کند.
\end{theorem}

این قضیه را در بخش \ref{sec:cyk} با الگوریتم \lr{CYK} ثابت می‌کنیم. با این حال $\abs{w}^3$ هم برای یک برنامه چند هزار خطی زیاد است. آرزو داریم الگوریتمی \textbf{خطی} (با زمان متناسب با $\abs{w}$) داشته باشیم. برای گرامرهای مستقل از متن در حالت کلی چنین روشی نمی‌شناسیم، اما برای حالت‌های خاص و مهمی وجود دارد.

\section{گرامرهای ساده}

\begin{definition}[گرامر ساده (\lr{s-grammar})]
گرامر مستقل از متن $G=(V,T,S,P)$ یک \textbf{گرامر ساده} است اگر همه قاعده‌هایش به شکل $A\to ax$ ($A\in V$، $a\in T$، $x\in V^*$) باشند و هر زوج $(A,a)$ حداکثر در یک قاعده ظاهر شود.
\end{definition}

\begin{example}
گرامر $S\to aS\mid bSS\mid c$ ساده است. اما $S\to aS\mid bSS\mid aSS\mid c$ ساده نیست، چون زوج $(S,a)$ در دو قاعده $S\to aS$ و $S\to aSS$ آمده است.
\end{example}

در گرامر ساده، تجزیه هر $w=a_1a_2\cdots a_n$ زمانی متناسب با $\abs{w}$ می‌گیرد. برای شروع تنها یک انتخاب داریم: قاعده‌ای به شکل $S\to a_1x$. پس اشتقاق با $S\Rightarrow a_1A_1\cdots A_m$ شروع می‌شود. برای جایگزینی $A_1$ هم تنها یک قاعده به شکل $A_1\to a_2x'$ وجود دارد و همین‌طور تا آخر. هر گام دقیقاً یک پایانه تولید می‌کند و هیچ انتخابی نداریم، پس کل کار در حداکثر $\abs{w}$ گام تمام می‌شود.

\section{الگوریتم \lr{CYK}}\label{sec:cyk}

الگوریتم \lr{CYK} (کوک، یانگر و کاسامی) یک الگوریتم عضویت و تجزیه با زمان $O(n^3)$ است. این الگوریتم فقط روی گرامرهای در \textbf{فرم نرمال چامسکی} کار می‌کند (که طبق قضیه \ref{thm:cnf} محدودیتی نیست) و یک مسئله را به دنباله‌ای از مسئله‌های کوچک‌تر می‌شکند (\textbf{برنامه‌نویسی پویا}).

فرض کنید $G=(V,T,S,P)$ در \lr{CNF} و $w=a_1a_2\cdots a_n$ باشد. زیررشته‌ها و مجموعه‌های زیر را تعریف می‌کنیم:
\[w_{ij}=a_i\cdots a_j,\qquad V_{ij}=\{A\in V\st A\der w_{ij}\}.\]
یعنی $V_{ij}$ همه متغیرهایی است که زیررشته $w_{ij}$ را تولید می‌کنند. پس
\[w\in L(G)\iff S\in V_{1n}.\]

\textbf{محاسبه $V_{ij}$:}
\begin{itemize}
\item \textbf{زیررشته‌های تک‌نمادی:} $A\in V_{ii}$ اگر و تنها اگر $G$ قاعده $A\to a_i$ را داشته باشد (چون در \lr{CNF} تنها راه تولید یک پایانه همین است).
\item \textbf{زیررشته‌های بلندتر} ($j>i$): $A$ رشته $w_{ij}$ را تولید می‌کند اگر و تنها اگر قاعده‌ای مثل $A\to BC$ داشته باشد و $w_{ij}$ را بتوان در نقطه‌ای مثل $k$ شکست که $B$ قسمت اول و $C$ قسمت دوم را تولید کند:
\[V_{ij}=\bigcup_{k\in\{i,\dots,j-1\}}\{A\st A\to BC,\ B\in V_{ik},\ C\in V_{k+1,j}\}.\]
\end{itemize}
این مجموعه‌ها را به ترتیب طول محاسبه می‌کنیم: اول $V_{11},V_{22},\dots,V_{nn}$، بعد $V_{12},V_{23},\dots,V_{n-1,n}$، بعد $V_{13},V_{24},\dots$ و الی آخر. وقتی $V_{ij}$ را حساب می‌کنیم، همه مجموعه‌های زیررشته‌های کوتاه‌تر آماده‌اند.

\textbf{درستی} با استقرا روی طول $j-i$ ثابت می‌شود: برای طول ۱ مستقیم از تعریف \lr{CNF}، و برای طول بیشتر، اولین گام هر اشتقاق از $A$ باید $A\Rightarrow BC$ باشد و $B$ و $C$ دو تکه پیوسته از $w_{ij}$ را تولید می‌کنند، که طبق فرض استقرا در $V_{ik}$ و $V_{k+1,j}$ ثبت شده‌اند.

اگر برای هر عضو $V_{ij}$ به خاطر بسپاریم از کدام $k$ و کدام قاعده آمده، الگوریتم \textbf{تجزیه} هم انجام می‌دهد: از $S\in V_{1n}$ به عقب برمی‌گردیم و درخت تجزیه را می‌سازیم.

\begin{example}
رشته $w=aabbb$ و گرامر
\[S\to AB,\qquad A\to BB\mid a,\qquad B\to AB\mid b.\]
\begin{itemize}
\item $V_{11}=V_{22}=\{A\}$ (قاعده $A\to a$) و $V_{33}=V_{44}=V_{55}=\{B\}$.
\item $V_{12}$: تنها شکست $a|a$ است و $AA$ سمت راست هیچ قاعده‌ای نیست، پس $V_{12}=\emptyset$.
\item $V_{23}$: شکست $a|b$ با $A\in V_{22}$ و $B\in V_{33}$؛ قاعده‌های با سمت راست $AB$: \ $S\to AB$ و $B\to AB$. پس $V_{23}=\{S,B\}$.
\item $V_{34}$: شکست $b|b$ با $BB$: \ $V_{34}=\{A\}$، و به همین شکل $V_{45}=\{A\}$.
\end{itemize}
با ادامه همین کار جدول کامل به دست می‌آید:
\begin{ltrfig}
\renewcommand{\arraystretch}{1.3}
\begin{tabular}{|c|c|c|c|c|}
\hline
$V_{11}=\{A\}$ & $V_{22}=\{A\}$ & $V_{33}=\{B\}$ & $V_{44}=\{B\}$ & $V_{55}=\{B\}$\\ \hline
$V_{12}=\emptyset$ & $V_{23}=\{S,B\}$ & $V_{34}=\{A\}$ & $V_{45}=\{A\}$ & \multicolumn{1}{c}{}\\ \cline{1-4}
$V_{13}=\{S,B\}$ & $V_{24}=\{A\}$ & $V_{35}=\{S,B\}$ & \multicolumn{2}{c}{}\\ \cline{1-3}
$V_{14}=\{A\}$ & $V_{25}=\{S,B\}$ & \multicolumn{3}{c}{}\\ \cline{1-2}
$V_{15}=\{S,B\}$ & \multicolumn{4}{c}{}\\ \cline{1-1}
\end{tabular}
\end{ltrfig}
مثلاً برای $V_{15}$ شکست $a|abbb$ را با $A\in V_{11}$ و $B\in V_{25}$ داریم که $S$ و $B$ را می‌دهد. چون $S\in V_{15}$، نتیجه می‌گیریم $aabbb\in L(G)$.
\end{example}

\textbf{تحلیل زمان:} دقیقاً $\frac{n(n+1)}{2}$ مجموعه $V_{ij}$ باید محاسبه شود و محاسبه هر کدام حداکثر $n$ نقطه شکست $k$ را بررسی می‌کند (برای هر شکست، کار با تعداد قاعده‌ها متناسب است که به $n$ بستگی ندارد). پس کل زمان $O(n^3)$ است.

\begin{summary}
\begin{itemize}
\item عضویت: آیا $w\in L(G)$؟ تجزیه: پیدا کردن اشتقاق $w$.
\item جست‌وجوی کامل ساده است اما نمایی است و بدون حذف قاعده‌های $\eps$ و واحد ممکن است متوقف نشود.
\item در گرامرهای ساده (\lr{s-grammar}) تجزیه در زمان خطی انجام می‌شود.
\item \lr{CYK} برای گرامرهای \lr{CNF}: $V_{ij}$ متغیرهای مولد $w_{ij}$ است، به ترتیب طول محاسبه می‌شود و $w\in L(G)\iff S\in V_{1n}$. زمان: $O(n^3)$.
\end{itemize}
\end{summary}
